\clearpage
\SourceLeaf{148}
\begin{center}
\textsc{Gauss} to \textsc{Schumacher}. Göttingen, 6 July 1840.
\end{center}

\noindent . . . . . . . . . . .

You know that I myself never set much store by the procedure I have used since 1794, which was later given the name Méthode des moindres quarrés. Understand me rightly. I am not speaking of the great service it renders, which is clear enough, but that is not how I appraise things. The reason, or the sense in which I attached little value to it, is that from the very beginning the idea struck me as so natural, so very obvious, that I never doubted in the least that many people who had had to deal with numerical calculation must have hit upon such a device on their own and used it, without thinking it worth their while to make much ado about so natural a thing. Above all it was \textsc{Tobias Mayer} who occurred to me, and I remember very distinctly that \textls{often}, when I spoke to others about my method (as I did, e.\,g.\ on many occasions during my student years, 1795--1798), I said that I would make the largest wager imaginable that \textsc{Tobias Mayer} had already used the same method in his calculations. I now know from those papers that I would have \textls{lost} the wager. They do indeed contain eliminations, e.\,g.\ of 3 unknown quantities from 4 or 5 equations, but carried out as the most commonplace reckoner would do them, without any trace of a subtler art.

\noindent . . . . . . . . . . . .

